%% paper.tex
%% Loss Aversion Shifts the Coordination Boundary in Cooperative MARL
%% All numbers verified against experimental results as of 2026-06-05.
%%
%% INCONSISTENCIES / REVIEWER NOTES (check before submission):
%%   [I1] Exp B-Reduced lambda=1 w_crit (0.586) differs from Exp A @300k w_crit (0.666)
%%        and Exp N null (0.685). This is intentional and explained in the text:
%%        Exp A / Exp N use alpha=beta=1 (linear); Exp B-Reduced uses alpha=beta=0.88.
%%        The difference IS the curvature effect. Make sure this is clear to the reader.
%%   [I2] Exp T 600k w_crit is unphysical (-0.60) due to narrow w-grid. We state
%%        w_crit < 0.50 as a bound and never report the -0.60 number.
%%   [I3] lambda=7 quasi-separation: we report bootstrap w_crit only; slope and p-value
%%        are suppressed. Clearly noted in the text.
%%   [I4] The curvature effect at lambda=1 in the decomposition is significant in the
%%        NEGATIVE direction (CI [-0.172, -0.033] excludes zero from above). The
%%        analyze_behavioral_excess.py "significant" column marks it "no" because it
%%        checks CI_lo > 0. The paper reports it as significant (beneficial).
%%   [I5] Figure file paths below assume compilation from the project root. Adjust
%%        \graphicspath or file paths before submitting to Overleaf.
%%   [I6] DiagnosticWrapper scripts node_1 at a fixed pull of 20 units. Effective
%%        learning agents are nodes 2-6 (5 agents). Paper describes this accurately.
%%   [I7] Total run count (1,124) does not include earlier pilot/phase experiments
%%        (Phases 0-5). Only the six main experiments are counted.
%%
%% FIXES APPLIED from cross-check v1 (2026-06-05):
%%   [M1] Metastability heatmap alpha,beta confound: caption and §R2 updated to note
%%        that the 600k row uses alpha=beta=0.88 (Exp T), while 150k/300k rows use
%%        alpha=beta=1 (Exp A). Exp 0b alpha=beta=1 numbers cited to bound the
%%        horizon-only effect. Table 1 Exp T description corrected.
%%   [m2] Single null baseline chosen: 0.685 (Exp N). "18.5 pp" corrected to "16.6 pp"
%%        (0.851 - 0.685 = 0.166), consistent with the decomposition table.
%%   [m3] "0.585" corrected to "0.586"; "+0.266" corrected to "+0.265".
%%   [m4] "0.016" corrected to "0.015" in appendix and §M2.
%%   [m5] \citep{terry2021pettingzoo} added in §E1.

\documentclass[11pt]{article}

\usepackage[margin=1in]{geometry}
\usepackage[T1]{fontenc}
\usepackage[utf8]{inputenc}
\usepackage{amsmath,amssymb,amsthm}
\usepackage{graphicx}
\usepackage{booktabs}
\usepackage{hyperref}
\usepackage{cleveref}
\usepackage[acronym]{glossaries}
\usepackage{csquotes}
\usepackage{microtype}
\usepackage{xcolor}
\usepackage{subcaption}
\usepackage[round]{natbib}

\hypersetup{
    colorlinks=true,
    linkcolor=blue,
    citecolor=blue,
    urlcolor=blue
}

% Notation shortcuts
\newcommand{\wcrit}{w_{\mathrm{crit}}}
\newcommand{\PT}{\mathrm{PT}}

\title{\textbf{Loss Aversion Shifts the Coordination Boundary in Cooperative MARL:\\
A Phase Diagram Analysis in a Multi-Echelon Supply Chain}}

\author{
Agriya Yadav\thanks{Ashoka University. Correspondence: \texttt{agriyayadav@gmail.com}.} \\
}

\date{\today}

\begin{document}

\maketitle

% -----------------------------------------------------------------------
\begin{abstract}
% X: what + why relevant
Cooperative multi-agent reinforcement learning (MARL) systems fail when
reward centralization falls below a critical threshold: agents converge to
a zero-order \emph{lazy-agent} policy, abandoning cooperation entirely.
Despite widespread recognition of this failure mode, the collapse boundary
has not been characterized as a quantitative, measurable object.
% Y: why hard
The threshold is not a fixed constant---it shifts with training horizon
(metastability) and, as we show, with agent-level prospect-theoretic loss
aversion ($\lambda$), making it insufficient to study reward structure or
behavioral bias in isolation.
% Z: our contribution
We construct the first quantitative phase diagram of cooperative MARL
coordination collapse in (reward centralization $w$, training horizon $T$,
loss aversion $\lambda$) space, using pre-registered collapse definitions
and bootstrap-estimated critical points with confidence intervals.
% 1a: key numbers
Across 1,124 controlled experiments on a multi-echelon supply chain
substrate, we identify $\wcrit = 0.666$ $[0.635, 0.694]$ at 300k steps
under no behavioral bias.
A three-way null-model decomposition reveals that probability weighting
alone ($\alpha{=}\beta{=}0.88$, $\lambda{=}1$) shifts $\wcrit$ \emph{left}
by $0.100$ $[0.033, 0.172]$---facilitating coordination---while loss
aversion at $\lambda{=}3$ and $\lambda{=}7$ shifts it \emph{right} by
$0.089$ $[0.041, 0.141]$ and $0.166$ $[0.135, 0.202]$ respectively.
Probability weighting and loss aversion exert statistically significant
effects in \emph{opposite directions}, a decomposition that, to our
knowledge, has not previously been demonstrated in cooperative MARL.
\end{abstract}

% -----------------------------------------------------------------------
\section{Introduction}
\label{sec:intro}

% X1: the problem and why it matters
Automated replenishment in multi-echelon supply chains is a canonical
cooperative MARL problem: agents must coordinate ordering decisions across
nodes with delayed feedback, partial observability, and misaligned local
incentives~\citep{oroojlooyjadid2022}.
Failures in such systems---agents that under-order, hoard, or
stop ordering altogether---generate the bullwhip effect, stockouts, and
cascading profit losses.
Deploying MARL agents that exhibit human-like behavioral biases, such as
prospect-theoretic loss aversion, is increasingly realistic given the
trend toward human-in-the-loop and behavior-aware AI
systems~\citep{lalmohammed2025,ethayarajh2024kto}.

% X2: the specific gap
A well-documented failure mode is \emph{coordination collapse}: agents
converge to a zero-action policy, abandoning cooperation
entirely~\citep{liu2023lazy,peng2023curo}.
While the phenomenon is recognized, no prior work has characterized the
collapse boundary as a quantitative, measurable object with confidence
intervals.
Existing work either treats the boundary qualitatively (``centralization
helps'') or studies algorithmic mitigation without locating the transition
point.

% Y: why it is hard to measure
The boundary is not a fixed scalar.
It depends on the reward structure ($w$, the fraction of global reward
shared across agents), the training horizon $T$, and---as we
demonstrate---on whether agents carry prospect-theoretic loss aversion
($\lambda$).
These factors interact: measuring any one in isolation does not
characterize the system.
Prospect theory decomposes into two components with distinct mathematical
forms~\citep{tversky1992}: a nonlinear probability weighting function
($\alpha, \beta$) and a loss aversion coefficient ($\lambda$).
Prior work applying prospect theory to cooperative MARL~\citep{lalmohammed2025}
integrates these components jointly into the algorithm, making it
impossible to attribute effects to either component separately.

% Z: our approach
We close this gap by constructing the first quantitative phase diagram of
cooperative MARL coordination collapse.
We treat the collapse boundary as a logistic curve in $w$-space, estimate
the critical point $\wcrit$ with bootstrap confidence intervals, and use
a pre-registered null-model design to decompose the total behavioral effect
into a \emph{probability weighting component} (curvature) and a
\emph{loss aversion increment} ($\lambda$-dependent).

% 1: what we find
Our main results are:
\begin{itemize}
  \item The structural coordination boundary is $\wcrit = 0.666$
        $[0.635, 0.694]$ at 300k steps under no behavioral bias, and
        $\wcrit = 0.822$ $[0.785, 0.858]$ at 150k steps, demonstrating
        that the boundary is horizon-dependent (metastable attractor).
  \item Probability weighting alone ($\alpha{=}\beta{=}0.88$,
        $\lambda{=}1$) \emph{reduces} $\wcrit$ by $0.100$
        $[0.033, 0.172]$ relative to the no-PT null, facilitating
        coordination.
  \item Loss aversion at $\lambda{=}3$ and $\lambda{=}7$ shifts
        $\wcrit$ by $+0.089$ $[0.041, 0.141]$ and $+0.166$
        $[0.135, 0.202]$ above the null, respectively.
        The non-overlapping confidence intervals of the $\lambda{=}7$
        and $\lambda{=}1$ conditions confirm the behavioral claim.
  \item The two PT components act in \emph{opposite directions}: curvature
        facilitates coordination; loss aversion impedes it.
\end{itemize}

% Figure 1 on page 1 (per Brendan's advice)
\Cref{fig:heatmap} shows the collapse probability surface in $(w,T)$
space, illustrating the metastability result.
\Cref{fig:decomposition} shows the behavioral decomposition: the four
logistic curves separate clearly, confirming both the curvature and
loss-aversion effects.

\begin{figure}[tbp]
  \centering
  \includegraphics[width=0.75\linewidth]{/Users/Agriya/Desktop/prospect theory/docs/expT_figures/collapse_heatmap.png}
  \caption{%
    $P(\mathrm{collapse})$ in $(w, T)$ space.
    Each cell shows the empirical collapse fraction across 20 seeds.
    The 150k and 300k rows are from Exp~A ($\alpha{=}\beta{=}1$); the
    600k row is from Exp~T ($\alpha{=}\beta{=}0.88$).
    Because $\alpha{=}\beta{=}0.88$ shifts $\wcrit$ left by 0.100
    (§\ref{sec:results_behavioral}), the apparent 300k$\to$600k slide
    is partially curvature-inflated.
    Under matched $\alpha{=}\beta{=}1$ conditions (Exp~0b),
    $P(\mathrm{collapse}\mid w{=}0.6) = 0.35$ and
    $P(\mathrm{collapse}\mid w{=}0.7) = 0.20$ at 600k,
    confirming that metastability holds with $\alpha{=}\beta$ held
    constant; only the magnitude is curvature-inflated in this figure.
  }
  \label{fig:heatmap}
\end{figure}

% -----------------------------------------------------------------------
\section{Background}
\label{sec:background}

% B1: MARL and CTDE
In cooperative MARL, $N$ agents share an environment and aim to maximize a
joint objective.
Centralised Training with Decentralised Execution
(CTDE)~\citep{lowe2017maddpg} addresses the \emph{credit assignment}
problem by allowing a centralised critic during training while agents
execute independently.
The degree of credit sharing is controlled by a reward mixing weight:
a scalar $w \in [0,1]$ that interpolates between purely local ($w{=}0$)
and purely global ($w{=}1$) reward.

% B2: lazy agents
When $w \approx 0$, each agent's reward is dominated by its own
performance.
The contribution of cooperative action is diluted: an agent receives no
signal when a teammate's action improves the global outcome.
\citet{liu2023lazy} and \citet{peng2023curo} formalise this as the
\emph{lazy agent} problem and \emph{relative overgeneralization}
respectively: agents converge to a safe, zero-action policy that avoids
coordination risk.
This paper asks: at what value of $w$ does this failure mode become
probable, and how does agent-level loss aversion shift that value?

% B3: prospect theory
Prospect theory~\citep{kahneman1979} models human decision-making under
uncertainty via a reference-dependent value function and a nonlinear
probability weighting function.
Cumulative prospect theory (CPT)~\citep{tversky1992} extends this to
compound lotteries via rank-dependent integration.
The CPT value function is:
\begin{equation}
  v(x) =
  \begin{cases}
    x^{\alpha} & \text{if } x \geq 0 \\
    -\lambda(-x)^{\beta} & \text{if } x < 0,
  \end{cases}
  \label{eq:cpt}
\end{equation}
where $\alpha, \beta \in (0,1]$ govern curvature (probability weighting
sensitivity) and $\lambda > 1$ is the loss aversion coefficient.
Setting $\alpha{=}\beta{=}1$ and $\lambda{=}1$ recovers linear expected
utility.
The parameters $\alpha, \beta$ and $\lambda$ have distinct effects on
agent behaviour: curvature introduces nonlinear sensitivity to outcome
magnitudes, while loss aversion introduces asymmetry between gains and
losses around the reference point.

% B4: CPT as reward transform
We apply CPT as a \emph{reward transform} applied after reward blending.
The per-step reward for agent $i$ is:
\begin{equation}
  r_i^{\PT} = v\!\left((1-w)\,r_i^{\mathrm{local}} + w\,r^{\mathrm{global}}\right),
  \label{eq:reward}
\end{equation}
where $v(\cdot)$ is given by \Cref{eq:cpt}.
This differs architecturally from~\citet{lalmohammed2025}, who integrate
CPT into the actor and critic update mechanics; our formulation applies PT
at the reward signal level and does not modify the policy gradient
computation.

% -----------------------------------------------------------------------
\section{Related Work}
\label{sec:related}

% R1: credit assignment / lazy agents -- closest structural siblings
The lazy agent problem~\citep{liu2023lazy} and relative
overgeneralization~\citep{peng2023curo} establish coordination collapse as
a canonical MARL pathology under CTDE, proposing algorithmic mitigations
(CMAE and CURO respectively).
\citet{weng2026} demonstrate that exploration noise structurally degrades
shared neural representations under decentralised training, explaining a
key mechanism of collapse.
\citet{dmarl2025} show empirically that even sophisticated decentralised
reward shaping cannot overcome fundamental coordination barriers, achieving
$-24.20$ average reward versus a centralised baseline at $+1.92$.
These works establish \emph{why} collapse occurs; none characterises
\emph{where} the collapse boundary lies as a quantitative object with
confidence intervals.
\citet{wang2025maca} construct counterfactual advantage functions at
multiple credit-assignment levels and evaluate robustness across
thresholding hyperparameters---a sweep that borders on a centralization
sensitivity analysis---but they optimise win rates on SMAC rather than
deriving a statistical phase boundary in $(w, T)$ space.
\citet{jiang2018dgn} and \citet{zhang2026lcerd} similarly utilise reward
mixing weights and observe phase-like transitions in training dynamics,
respectively, but neither isolates a named critical centralization value
with confidence intervals.

% R2: phase structure in MARL -- closest to ours, different axis
Most recently, \citet{yamaguchi2025} construct a phase diagram for
independent Q-learning across environment size $L$ and agent density
$\rho$, identifying coordinated, fragile, and jammed regimes separated by
an instability ridge.
Our work is orthogonal: we hold the environment and algorithm fixed and
map the phase boundary as a function of the reward structure parameter $w$
and the behavioral parameter $\lambda$---dimensions their framework does
not address.

% R3: PT in cooperative MARL -- Lalmohammed
\citet{lalmohammed2025} integrate CPT into the actor and critic update
mechanics of MADDPG, demonstrating that tunable risk profiles are
achievable in cooperative settings.
Our work is architecturally distinct (CPT as reward transform, not
gradient transform) and asks a different question: not whether PT-equipped
agents can coordinate, but at what centralization level coordination fails
and how the PT parameter $\lambda$ shifts that level.

% R4: risk-sensitive MARL -- CVaR approaches
\citet{qiu2021rmix} and \citet{xie2023riskq} incorporate Conditional Value
at Risk (CVaR) into cooperative value factorisation.
CVaR optimises the expected return in the worst $\alpha$-percentile of
outcomes; loss aversion operates via reference-dependent utility
asymmetry.
These are mathematically distinct mechanisms: CVaR addresses aleatoric
tail risk, while $\lambda$ captures the perceived asymmetry between gains
and losses~\citep{xie2023riskq}.

% R5: behavioral collapse in adjacent fields
\citet{mengesha2026} prove in evolutionary game theory that loss aversion
induces population-level lock-in at zero-sum equilibria, establishing the
conceptual phenomenon in a theoretical setting.
We extend this to deep MARL with learned neural policies and continuous
state spaces, providing the first empirical measurement in this setting.

% R6: supply chain MARL
\citet{oroojlooyjadid2022} establish the Beer Distribution Game as a
multi-agent Dec-POMDP where decentralised DQN triggers systemic cost
escalation corrected by a shared penalty.
\citet{gijsbrechts2022} demonstrate that A3C can match state-of-the-art
heuristics in multi-echelon inventory settings, establishing DRL as a
viable approach for supply chain optimisation.
\citet{marlim2023} apply MARL to multi-product single-echelon inventory.
Our substrate generalises these to a six-node multi-echelon network.

% -----------------------------------------------------------------------
\section{Method}
\label{sec:method}

% M1: the w parameter
The reward centralization parameter $w \in [0,1]$ is the sole structural
intervention.
All other environment parameters, algorithm hyperparameters, and random
seeds are held constant across conditions.

% M2: collapse definition (pre-registered)
We define coordination collapse as a binary outcome:
\begin{equation}
  \mathrm{collapsed}(\text{run}) = \mathbf{1}\!\left[
    \bar{S} < 10 \;\wedge\; \bar{\pi} \leq -127.1
  \right],
  \label{eq:collapse}
\end{equation}
where $\bar{S}$ is the mean base-stock level and $\bar{\pi}$ is mean
profit over ten evaluation episodes.
This definition was pre-registered before any analysis (commit
\texttt{03c3fc6}).
\Cref{sec:appendix_robustness} demonstrates robustness to $\pm 20\%$
threshold perturbation (maximum $\Delta\wcrit = 0.015$).

% M3: logistic phase boundary estimation
For each experimental condition, we fit:
\begin{equation}
  P(\mathrm{collapsed} \mid w) = \sigma(b_0 + b_1 w),
  \label{eq:logistic}
\end{equation}
where $\sigma$ is the logistic function.
We estimate $\wcrit = -b_0/b_1$ (the value of $w$ where collapse
probability equals 0.5) and its 95\% confidence interval via 2,000
bootstrap resamples at the run level.
We report Wilson score 95\% confidence intervals for per-$w$ collapse
fractions.

% M4: three-way decomposition -- the key methodological contribution
To attribute the observed shift in $\wcrit$ to specific PT components, we
define three conditions with identical environments and algorithms:
\begin{enumerate}
  \item \textbf{Null (Exp~N):} $\alpha{=}\beta{=}1$, $\lambda{=}1$---the
        linear, no-PT condition.
  \item \textbf{Curvature (Exp~B, $\lambda{=}1$):}
        $\alpha{=}\beta{=}0.88$, $\lambda{=}1$---probability weighting
        with no loss aversion.
  \item \textbf{Loss aversion (Exp~B, $\lambda \in \{3,7\}$):}
        $\alpha{=}\beta{=}0.88$, $\lambda > 1$---full PT.
\end{enumerate}
The curvature effect is defined as
$\wcrit(\text{B},\lambda{=}1) - \wcrit(\text{N})$,
and the loss aversion increment as
$\wcrit(\text{B},\lambda) - \wcrit(\text{B},\lambda{=}1)$.
This decomposition is valid because the null removes all PT distortion,
while the curvature condition adds nonlinearity without loss asymmetry.

% -----------------------------------------------------------------------
\section{Experimental Setup}
\label{sec:setup}

% E1: environment
We use \texttt{NetworkManagement-v1} from or-gym~\citep{hubbs2020orgym}:
a six-node multi-echelon inventory network with Poisson demand
($\mu{=}20$) on the retail edge, 30-period episodes, and uniform lead
times.
Node~1 (the retail-facing node) follows a fixed demand-pull policy,
stabilising the demand signal; nodes~2--6 are trained as cooperative PPO
agents.
Actions are parameterised via a base-stock wrapper that maps each
agent's continuous action to a target inventory level $S \in [0, 500]$,
with order quantity $Q = \max(0, S - \mathrm{IP})$ where $\mathrm{IP}$
is inventory position.

% E2: algorithm
Each agent runs PPO~\citep{schulman2017ppo} with a shared MLP policy
(\texttt{n\_steps=128}, \texttt{batch\_size=256}) via
Stable-Baselines3~\citep{raffin2021sb3}, PettingZoo~\citep{terry2021pettingzoo},
and SuperSuit~\citep{terry2020supersuit} vectorisation.
Evaluation is deterministic over 10 episodes with 20 seeds per condition.

% E3: experiment inventory
\Cref{tab:experiments} summarises all six experiments.
All designs were pre-registered before execution.
Total run count: 1,124.

\begin{table}[t]
\centering
\caption{Experiment summary. All experiments use the same environment and
algorithm.  $\alpha{=}\beta{=}1$ denotes the linear (no-PT) condition.}
\label{tab:experiments}
\small
\begin{tabular}{llrrrrr}
\toprule
Exp & Purpose & $w$ values & $\lambda$ & $\alpha,\beta$ & Seeds & Runs \\
\midrule
A @150k  & Structural boundary  & $\{0.0, 0.1,\ldots, 1.0\}$ & 1 & 1,1 & 20 & 220 \\
A @300k  & Structural boundary  & $\{0.2,\ldots, 0.8\}$       & 1 & 1,1 & 20 & 140 \\
C        & Basin control        & $\{0.5, 0.85, 1.0\}$        & 1 & 1,1 & 12 & 144 \\
N        & Null baseline        & $\{0.0, 0.2,\ldots, 1.0\}$  & 1 & 1,1 & 20 & 180 \\
B-Red.   & Behavioral surface   & $\{0.5, 0.65, 0.75, 0.85, 0.95\}$ & $\{1,3,7\}$ & 0.88 & 20 & 300 \\
T (600k only) & Metastability horizon & $\{0.5, 0.55,\ldots, 0.80\}$ & 1 & 0.88 & 20 & 140 \\
\bottomrule
\end{tabular}
\end{table}

% -----------------------------------------------------------------------
\section{Results}
\label{sec:results}

\subsection{Structural Phase Transition and Metastability}
\label{sec:results_structural}

% R1: logistic transition is real and sharp
Exp~A establishes a sharp logistic collapse transition.
At 150k steps, $\wcrit = 0.822$ $[0.785, 0.858]$ with transition width
$\Delta w = 0.278$ (logistic slope $b_1 = -15.78$, $p = 5 \times 10^{-9}$).
At 300k steps, $\wcrit = 0.666$ $[0.635, 0.694]$ with $\Delta w = 0.199$
($b_1 = -22.05$, $p = 3.2 \times 10^{-7}$).
In both cases, collapse fractions rise steeply from near-zero to near-one
within a $\Delta w < 0.3$ range, consistent with a genuine phase
transition rather than a gradual degradation.

% R2: metastability -- the boundary slides with T
The critical boundary decreases monotonically with training horizon $T$:
$\wcrit = 0.822$ at 150k and $\wcrit = 0.666$ at 300k (both from Exp~A,
$\alpha{=}\beta{=}1$).
\Cref{fig:heatmap} visualises the collapse surface in $(w, T)$ space.
This constitutes a \emph{metastability} result: the lazy-agent attractor
is not absorbing.
Agents trapped below the isocline at finite $T$ can escape given
sufficient training; $w$ controls the energy barrier height of the
attractor.
The sharpening transition width ($\Delta w$ decreasing from 0.278 to
0.199 across 150k to 300k) is consistent with policy convergence reducing
outcome variance.

\paragraph{Note on the 600k row.}
The 600k row in \Cref{fig:heatmap} comes from Exp~T, which uses
$\alpha{=}\beta{=}0.88$.
Because our decomposition shows $\alpha{=}\beta{=}0.88$ shifts
$\wcrit$ left by 0.100 relative to $\alpha{=}\beta{=}1$, the 600k
collapse fractions in the figure are partially curvature-deflated, not
purely a horizon effect.
To bound the horizon-only slide, we use Exp~0b ($\alpha{=}\beta{=}1$,
20~seeds, 600k): $P(\mathrm{collapse}\mid w{=}0.6) = 0.35$ and
$P(\mathrm{collapse}\mid w{=}0.7) = 0.20$, compared to 0.80 and 0.40
at 300k.
The metastability result---the boundary slides left with $T$---holds
under matched $\alpha{=}\beta{=}1$ conditions; only the magnitude at
600k in the figure is inflated by the curvature effect.

% R3: basin control rules out initialization
Exp~C rules out initialization as a driver of collapse.
Within-$w$ correlations between $S_{\mathrm{init}}$ and converged $S$
are $-0.11$ ($w{=}0.5$), $-0.38$ ($w{=}0.85$), and $-0.36$ ($w{=}1.0$)---
all near zero and inconsistent with initialization dominating the outcome.
The collapsed/survived assignment is not sensitive to where agents start.

\subsection{Null Model and Reproducibility}
\label{sec:results_null}

% R4: null model reproduces
Exp~N (null: $\alpha{=}\beta{=}1$, $\lambda{=}1$) yields
$\wcrit = 0.685$ $[0.657, 0.710]$, compared to $\wcrit = 0.666$
$[0.635, 0.694]$ for the PT ($\lambda{=}1$, $\alpha{=}\beta{=}1$) condition
from Exp~A @300k---a difference of $-0.020$ with overlapping confidence
intervals.
Both conditions use no loss aversion and linear value functions.
Their agreement to within 0.020 confirms that the $\wcrit$ measurement
reproduces across independent experimental runs, validating the null
baseline.

\subsection{Behavioral Phase Boundary}
\label{sec:results_behavioral}

% R5: behavioral claim confirmed
Exp~B-Reduced yields monotonically increasing $\wcrit$ with loss aversion
$\lambda$:
$\wcrit(\lambda{=}1) = 0.586$ $[0.539, 0.624]$,
$\wcrit(\lambda{=}3) = 0.774$ $[0.751, 0.798]$, and
$\wcrit(\lambda{=}7) = 0.851$ $[0.845, 0.858]$.\footnote{%
  For $\lambda{=}7$, the near-complete separation of outcomes
  (collapse fractions 1.00, 1.00, 1.00, 0.55, 0.00 at
  $w \in \{0.5, 0.65, 0.75, 0.85, 0.95\}$) renders the logistic slope
  and p-value unreliable.
  We report $\wcrit$ from bootstrap estimation, which is robust to
  quasi-separation; the slope and p-value are not reported for this
  condition.}
The behavioral claim is confirmed: $\wcrit(\lambda{=}7)$ lower CI
$0.845$ exceeds $\wcrit(\lambda{=}1)$ upper CI $0.624$---the confidence
intervals are non-overlapping by 0.221.
\Cref{fig:wcrit_lambda} illustrates the monotone relationship.

\begin{figure}[tbp]
  \centering
  \includegraphics[width=0.82\linewidth]{/Users/Agriya/Desktop/prospect theory/docs/expB_reduced_figures/fig2_behavioral_decomposition.png}
  \caption{%
    Behavioral decomposition of the collapse boundary.
    Each curve shows $P(\mathrm{collapse} \mid w)$ at 300k steps.
    The null ($\alpha{=}\beta{=}1$, gray dashed) is the structural
    baseline.
    The $\lambda{=}1$, $\alpha{=}\beta{=}0.88$ curve (blue) shifts
    \emph{left} of the null by 0.100, indicating that probability
    weighting alone facilitates coordination.
    The $\lambda{=}3$ (orange) and $\lambda{=}7$ (red) curves shift
    progressively \emph{right}, indicating that loss aversion raises the
    centralization requirement.
    Vertical dashed lines mark $\wcrit$ estimates; shaded bands show
    bootstrap 95\% confidence intervals.
  }
  \label{fig:decomposition}
\end{figure}

\begin{figure}[tbp]
  \centering
  \includegraphics[width=0.6\linewidth]{/Users/Agriya/Desktop/prospect theory/docs/expB_reduced_figures/expB_wcrit_vs_lambda.png}
  \caption{%
    $\wcrit$ as a function of loss aversion $\lambda$ (Exp~B-Reduced,
    $\alpha{=}\beta{=}0.88$, 300k steps, 20 seeds per cell).
    Error bars show bootstrap 95\% confidence intervals.
    The monotone increase and non-overlapping CIs between $\lambda{=}1$
    and $\lambda{=}7$ confirm the behavioral claim.
    Note: the $\lambda{=}7$ CI $[0.845, 0.858]$ should be interpreted
    with caution due to quasi-complete separation; see
    \Cref{sec:results_behavioral}.
  }
  \label{fig:wcrit_lambda}
\end{figure}

% R6: three-way decomposition
\Cref{tab:decomposition} and \Cref{fig:decomposition} present the
three-way decomposition.
Relative to the null ($\wcrit = 0.685$):

\begin{enumerate}
  \item \textbf{Curvature effect} ($\alpha{=}\beta{=}0.88$, $\lambda{=}1$):
        excess $= -0.100$, 95\% CI $[-0.172, -0.033]$.
        The confidence interval excludes zero from above, indicating that
        probability weighting alone \emph{significantly reduces} the
        required centralization.
        Agents with nonlinear probability perception but no loss aversion
        are \emph{easier} to coordinate than agents with linear utility.

  \item \textbf{Loss aversion at $\lambda{=}3$}:
        excess $= +0.089$, 95\% CI $[+0.041, +0.141]$
        (CI excludes zero; significant).

  \item \textbf{Loss aversion at $\lambda{=}7$}:
        excess $= +0.166$, 95\% CI $[+0.135, +0.202]$
        (CI excludes zero; significant).
\end{enumerate}

The curvature and loss-aversion effects operate in \emph{opposite
directions}: curvature facilitates coordination, loss aversion impedes it.
At $\lambda{=}7$, the loss aversion increment above the curvature baseline
is $+0.265$ ($= 0.851 - 0.586$), overwhelming the curvature benefit
of $-0.100$ for a net shift of $+0.166$ above the structural null
(rounding: $-0.100 + 0.265 \approx 0.165 \approx 0.166$).

\begin{table}[t]
\centering
\caption{Three-way behavioral decomposition of $\wcrit$ relative to the
null model. ``Significant'' denotes that the conservative difference
confidence interval excludes zero. $^{\dagger}$For $\lambda{=}7$,
$\wcrit$ is from bootstrap only; see footnote in
\Cref{sec:results_behavioral}.}
\label{tab:decomposition}
\small
\begin{tabular}{lccccc}
\toprule
Condition & $\wcrit$ & 95\% CI & Excess & Excess CI & Sig.? \\
\midrule
Null ($\alpha{=}\beta{=}1$, $\lambda{=}1$) & 0.685 & [0.657, 0.710] & 0 & --- & --- \\
$\alpha{=}\beta{=}0.88$, $\lambda{=}1$ & 0.586 & [0.539, 0.624] & $-0.100$ & $[-0.172, -0.033]$ & $\checkmark$ (neg.) \\
$\alpha{=}\beta{=}0.88$, $\lambda{=}3$ & 0.774 & [0.751, 0.798] & $+0.089$ & $[+0.041, +0.141]$ & $\checkmark$ \\
$\alpha{=}\beta{=}0.88$, $\lambda{=}7$$^{\dagger}$ & 0.851 & [0.845, 0.858] & $+0.166$ & $[+0.135, +0.202]$ & $\checkmark$ \\
\bottomrule
\end{tabular}
\end{table}

% -----------------------------------------------------------------------
\section{Discussion}
\label{sec:discussion}

% D1: the opposite-effects finding is the main result
The finding that probability weighting and loss aversion exert opposite
effects on the coordination boundary is, to our knowledge, novel.
Prior work applies prospect theory as a unified behavioural model without
decomposing its components~\citep{lalmohammed2025,anon2024risk}.
Our null-model design isolates the two components, revealing that the
beneficial effect of curvature ($-0.100$) is masked when combined with
loss aversion ($\lambda{>}1$), which reverses the sign and amplifies the
required centralization.
This has a concrete operational interpretation: a system with $\lambda{=}7$
agents---calibrated to reflect human loss aversion
estimates~\citep{tversky1992}---requires $w > 0.851$ for reliable
coordination, compared to $w > 0.666$ for risk-neutral agents.
The additional 16.6 percentage points of centralization attributable to
behavioral bias ($\wcrit(\lambda{=}7) - \wcrit(\text{null}) =
0.851 - 0.685 = 0.166$, consistent with the excess in
\Cref{tab:decomposition}) are a quantified cost of deploying loss-averse
agents in a cooperative system.

% D2: metastability -- lazy agent attractor is escapable
The sliding $\wcrit$ result is a positive finding, not a pathology.
It establishes that the lazy-agent attractor is \emph{metastable}: agents
caught below the isocline at finite $T$ can, in principle, coordinate if
given more training.
This is consistent with \citet{weng2026}'s finding that exploration breaks
cooperation in shared-policy MARL during early training, before
representations converge.
The practical implication is that the minimum $w$ required depends on the
operational training budget: tighter horizons demand higher centralization.

% D3: relation to Yamaguchi 2025
\citet{yamaguchi2025} construct a phase diagram in $(L, \rho)$ space for
IQL, identifying analogous coordinated/fragile/jammed regimes.
Our results complement theirs by demonstrating that the phase structure is
not only a function of environmental complexity but also of the reward
structure ($w$) and agent behavioral parameters ($\lambda$).
Together, these works suggest that MARL coordination exhibits coherent
phase structure across multiple parameter axes.

% D4: limitations
Several limitations should be noted.
First, all results are demonstrated on a single environment
(NetworkManagement-v1) and algorithm (PPO); replication across environments
and algorithms remains future work.
Second, the $\lambda{=}7$ condition exhibits quasi-complete separation;
additional $w$ values above 0.85 would sharpen the $\wcrit$ estimate.
Third, the metastability result describes finite-horizon behavior; an
asymptotic lower bound on $\wcrit$ is not estimated by our experiments.
Fourth, we fix the demand distribution and lead times; the dependence of
$\wcrit$ on these environmental parameters is unexplored.

% -----------------------------------------------------------------------
\section{Conclusion}
\label{sec:conclusion}

We presented the first quantitative phase diagram of cooperative MARL
coordination collapse, characterising the critical centralization threshold
$\wcrit$ as a function of reward structure, training horizon, and
prospect-theoretic loss aversion.
A three-way null-model decomposition reveals that the two components of
prospect theory---probability weighting and loss aversion---exert
statistically significant effects in opposite directions on the
coordination boundary: the former facilitates coordination by $0.100$;
the latter impedes it by up to $0.166$.
These results provide concrete design constraints for deploying cooperative
MARL in behavioral settings and establish a pre-registered phase-diagram
methodology applicable to any cooperative MARL system.

Natural extensions include replication across algorithms (MAPPO, QMIX)
and environments (MPE, SMAC), a mean-field theoretical derivation of the
logistic boundary shape, adaptive centralization scheduling based on the
estimated isocline, and characterization of heterogeneous-$\lambda$
populations.

% -----------------------------------------------------------------------
% \section*{Acknowledgements}
% \label{sec:ack}
% [To be completed.]

% -----------------------------------------------------------------------

% -----------------------------------------------------------------------
\appendix

\section{Threshold Robustness}
\label{sec:appendix_robustness}

We verify that the collapse definition (\Cref{eq:collapse}) is robust to
perturbation of its thresholds.
We re-estimate $\wcrit$ on Exp~A @300k data under nine variants of the
definition, varying $\bar{S}$ threshold over $\{8, 10, 12\}$ and
$\bar{\pi}$ threshold over $\{-102.5, -127.1, -152.5\}$.
Three variants with threshold $\bar{\pi} \leq -152.5$ are degenerate
(the minimum profit in the dataset is $-128.1$, so no run satisfies
this condition); they are excluded.
Across the remaining six estimable variants, the maximum deviation from
the baseline $\wcrit = 0.666$ is $0.015$, well below the confidence
interval width.
We conclude that the collapse definition is robust to $\pm 20\%$ threshold
perturbation.

% -----------------------------------------------------------------------
\bibliographystyle{plainnat}
\bibliography{references}

\end{document}